% Version 3.10: full executable ReadyBundle plus one prepared opening certificate.
% This fragment does not redefine def:ready or alg:bootstrap.
\paragraph{A complete regional ReadyBundle.}
\label{sec:ready-bundle}
The relay used in \cref{sec:ready} transfers the complete executable
state and one prepared opening certificate. It is not a READY quorum
or public activation certificate and does not change the public opening
of an export duty. Its timing claim has the complete-message and native
validation scope stated below.

For target session $(\beta,r)$, let $A=A_{\beta,r,0,0}$ and
$\Delta=\Delta_{i_{\beta,r}}$. A \tagg{READYBUNDLE} is an authenticated
regional message with the following typed fields. Every required object
is present as bytes, once, in its normalized dependency set. A URL,
state root, unresolved hash or promise to fetch later is not a substitute.
\begin{table}[htbp]\centering\small
\begin{tabularx}{\linewidth}{@{}>{\raggedright\arraybackslash}p{.25\linewidth}Y@{}}\toprule
Field & Required value\\\midrule
Domain and target & Chain and protocol domain, encoding/native-state
format identifier, target $(\beta,r)$, selected region and sender identity.\\
Fixed handoff & The complete handoff of \cref{sec:witness}, including
its cause, predecessor/source context, selection and fixed target.\\
Exact anchor & $A$'s original producing pair, height, hash, full DATA,
original direct or covering commitment evidence and required ancestry.\\
Executable state $S_A$ & The full, already constructed application-state
object from which execution after $A$ can continue. It includes all state,
code and continuation metadata required by the native execution interface;
if partial execution is supported, its exact continuation position is
included. It contains no imported regional locks, votes or entry timers.\\
State evidence $\Pi_A$ & Native validation evidence connecting this exact
$S_A$ to this exact $A$, with all proof inputs and dependencies. Use the
existing sound application-validation interface; a sender's assertion
or a bare state hash is insufficient.\\
Authority and bounds & Finalized $\tagg{AUTH}_{r,0}$, fixed configuration
and quota, $Q^0_{\beta,r}$, $q^{\rm open}_{\beta,r}$ and
$\lambda^{\rm open}_{\beta,r}$, and all authority evidence needed to
validate the inherited history.\\
Public evidence & An authorized replay checkpoint and complete finalized
L1 control/authority and Espresso header/namespace histories needed to
verify the handoff, references and supplied authorization context. The
checkpoint's trust/authorization dependencies are included. The packet contains the finalized opening header and its evidence.
It need not contain a header after the opening; the successor selection
floor advances only after the current terminal result forms.\\
Prepared opening input & One complete $\mathit{OC}$ satisfying
$\tagg{CheckOC}(A,\tagg{AUTH}_{r,0},\mathit{OC})$ for the exact preparation
key $K$ of \cref{sec:input-selection}; it contains the signed offer keys
and one selected authenticated request body if nonempty. It contains no
pending-pool pointer, private selection job, second hydration signature
set or historical request-owner directory.\\
Dependency set & Full bytes for the preceding fields' transitive
validation dependencies. A receiver need not fetch an object, replay
application execution to reconstruct $S_A$, or obtain a new authority
before processing this bundle.\\\bottomrule
\end{tabularx}
\end{table}

The native interface determines how $\Pi_A$ verifies state at $A$.
If the block format does not directly commit to a state root, the
bundle supplies the execution evidence and authenticated starting state
required by that interface; this contract does not invent such a
commitment. A merely serialized snapshot that still needs restoration
or conversion is not yet the directly executable state required here.
Recipients independently validate the packet. They do not trust the
sender to authorize a handoff or assert state correctness.

\paragraph{Send and receive events.}
Each correct selected member holds its complete bundle and a matching
valid opening certificate before its first Ready/Enter event for this
pair. Gathering histories, cross-region fetching, state reconstruction,
hydration, finite selection and bundle formation precede that event;
the existing model assigns them no fixed time bound. The pre-Ready
\tagg{STATE\_CANDIDATE} path of \cref{sec:input-selection} distributes
the executable state without waiting for this certificate or Ready.
At the event it
installs its verified executable state, attaches the validated certificate,
executes the existing initial Enter once, and actually sends the complete
bundle to every member of the fixed selected-region snapshot.
Every correct member does this at its own
first event, so the first correct entrant necessarily sends without
having to know that it is first. Bundle construction or asynchronous
dispatch is not deferred until after the recorded entry time.

The receiver processes the bundle through its serialized Session Manager:
\begin{enumerate}[label=(\roman*),itemsep=2pt]
\item Authenticate the domain, sender and target, and inspect supplied
evidence under the existing evidence-first rule. Validate the historical
handoff, exact anchor, all required authorities, references and native
state evidence. Check the matching $K$ and opening certificate under
\tagg{CheckOC}; no private pool, queue minimum or unfinished selection
job is consulted. Incomplete or invalid material does not authorize entry.
An existing compatible cache can eliminate duplicate validation, but the
sender does not rely on an unknown receiver cache to omit required bytes.
\item Apply the existing live-activation check against the receiver's
processed public state. A closed or superseded target, a conflicting fixed
handoff, a close request or an exit record prevents initial Enter. This
packet cannot replace a fixed anchor or override newer verified closure.
If the target is closed, invoke the manager's phase-independent closure
relay using this complete packet's verified historical snapshot before
refusing entry. This also applies at a Terminated record and after restart;
the worker need not be recreated. If the same pair already entered,
retain useful evidence without resetting
its state or entry time.
\item For an unentered valid target, validate and adopt the received
$S_A$, attach its complete validated opening certificate, check every
remaining condition of \cref{def:ready}, then execute
the existing Ready/Enter step of \cref{alg:bootstrap} in this same
message-processing event. The durable entry flag, fixed state attachment
and entry time are recorded atomically before signatures are dispatched.
Relay the complete bundle as above. A member need not finish its own
selection or independent reconstruction job after the supplied executable
object passes the original native predicate. No separate Ready queue,
state-fetch task or post-Ready installation job lies between this event
and Enter. Late callbacks cannot replace the attachment or reset entry
time; redundant jobs may be cancelled while all storage and promised
service/relay duties remain owned.
\end{enumerate}

The attachment is a cache for later sends, not a consensus choice of
one exclusive certificate. A valid later NEWVIEW/proposal can supply
another matching $\mathit{OC}$; validators cannot veto it merely because
their local attachment differs. Selected bytes may additionally enter the
canonical pending service, but REPORT does not wait for private queue
insertion or reindexing. A stale certificate remains historical evidence;
it never overrides activation or closure checks.

Here \emph{adopt} has a specific logical meaning: the validated received
state object becomes this new program's executable application state.
The program operates on that object; it does not first reconstruct a
different database from it. This is a representation and message-handler
contract, not a claim that arbitrary database import, disk copying or
native-runtime initialization takes zero time. No receiver can pass Ready
while required installation work remains unfinished.

\paragraph{The exact timing implication.}
Let $e$ be the first correct initial Enter time for this pair. The send
rule places a complete correct-to-correct bundle send at $e$. Suppose the
fixed handoff remains activatable throughout the interval in question
and no closure or other valid entry prohibition intervenes.
If validation, adoption and the atomic Enter transition are the complete
message-processing operation covered by \cref{ass:network}, then
\[
 e_v(\beta,r,0,0)\le e+\Delta
 \qquad\text{for every correct selected member }v.
\]
Indeed, the first correct sender supplies the full executable object and
all validation dependencies, including one opening certificate and at
most one selected request. By the existing delivery-and-processing bound
for this admitted concrete complete packet, each correct recipient
finishes that handler by $e+\Delta$.
Its guarded final step is Enter, or it has already entered earlier.
There is no remaining private-pool selection or deferred state-restoration
action in this handler argument. Its concrete resource scope is stated
below.

This is the supported abstract-message interpretation, not an additional
service bound for state installation. The quoted network premise covers
delivery and processing of complete validation dependencies. It does not,
by itself, bound an external native import or storage operation merely
because the protocol calls that operation ``installation.'' If such work
remains after bounded message validation, let $u_v\ge0$ be its actual
residual delay. The available decomposition is only
\[
 e_v(\beta,r,0,0)\le e+\Delta+u_v,
\]
when entry eventually completes and remains authorized; the frozen model
gives no uniform bound on $u_v$. Calling that work one event cannot prove
$u_v=0$. In that implementation, the $e+\Delta$ claim and the bootstrap
bounds derived from it remain unsupported.

Consequently the message contract repairs the missing protocol interface
at the abstract level only when the native state is genuinely an executable
validation dependency adopted by the modeled handler. A concrete native
adapter must demonstrate that it meets this contract; sound state proofs
alone do not demonstrate a timing property. The document must withhold
an unconditional native-installation claim until that obligation is met.
No claim is made about an arbitrary history catch-up outside the complete
packet or about when any member first constructs such a packet.

ReadyBundle is regional transport of existing state and evidence. It does
not add a quorum of READY signatures, an Espresso transaction, a publisher
guarantee or a public activation certificate. It neither starts nor resets
an export window. The public clock and the separate live-activation checks
remain those of \cref{alg:force,alg:external,sec:ready}.


\paragraph{Actual input and processing scope.}
\label{sec:input-resource-scope}
Let $n$ be the target roster size. An opening certificate supplies one
$K$, at most $n$ member signatures and signed keys, a distinct-signer
representation, and at most one canonical request body. A concrete
canonical framing uses the typed $K$ once, an $n$-bit signer bitmap,
entries in authenticated roster order, and a length-prefixed selected
body when nonempty. Each entry contains an EMPTY/nonempty tag, canonical
length and digest when nonempty, and that member's ordinary signature.
Every signed preimage includes its type, domain, $K$ and entry. The
bitmap and count must agree; no signature aggregation is assumed.
With $h_0$ the hash width, $s_v$ the ordinary encoded signature length,
$\ell(L)$ the canonical encoded-length size, and $c_{\rm frame}$ the
framing size, this gives
\[
\begin{split}
 \operatorname{bytes}(\mathit{OC})={}&c_{\rm frame}+c(K)+\lceil n/8\rceil\\
 &+\sum_{v\in S}\left[s_v+1+mathbf 1_{k_v\ne\tagg{EMPTY}}
                 (\ell(L_v)+h_0)\right]\\
 &+\mathbf 1_{k(\mathit{OC})\ne\tagg{EMPTY}}
                 (\ell(L_{\rm selected})+L_{\rm selected}).
\end{split}
\]
References into a normalized dependency set count their actual encoded
indices and require their full referents in the same packet. The formula
is for this explicit framing, not an assumption that omitted dependencies
are free. A singleton DATA fit does not prove that a maximal original
ReadyBundle plus that request and the signatures fits an old envelope.

The added Ready handler first parses the packet and binds $K$ to the
validated anchor, authority and native domain. It checks at most $n$
distinct signatures and weights and takes the minimum of the supplied
keys. It checks the selected body's authentication, canonicalization,
admission, fit and ancestry exclusion, then durably attaches the witness.
It does not scan a journal, search a live pool, verify a directory or
reconstruct a signer's cut. Original state validation and adoption remain
real work under their existing interface; no new native restoration is
part of this added handler.

An opening NEWVIEW carries one certificate and one initializer
advertisement signature; all recipients check that same choice.
An opening REPORT carries $\mathit{RC}$, a mode, certificate digest and
selected key, plus one report signature. A descendant REPORT carries
$\mathit{RC}$, \tagg{DISCHARGED}, \tagg{EMPTY} and its signature. A correct
leader checks at most $n$ reports against its one advertised certificate.
The first proposal supplies one certificate, one advertisement and one
REPORT signature set; its selected request is already the first batch
entry and is encoded once. Complete descendant/recovery histories retain
the first-DATA sidecars once per normalized dependency set. No path checks
one different certificate for every reporter. The added active validation
footprint is at most $n$ offer signatures, one advertisement signature,
and $n$ report signatures, together with the one selected-body check and
ordinary batch checks.

For fixed advertisement framing $c_{\rm ad}$ its bytes are
$c_{\rm ad}+\operatorname{bytes}(\mathit{RC})+h_0+
\operatorname{bytes}(k(\mathit{OC}))+s_{L(V)}$.
The opening REPORT amendment adds its mode and digest to the exact
context/key/signature framing. These are actual bytes and checks.
For an implementation with the original packet limits
$B_{\rm Ready},B_{\rm Newview},B_{\rm Report},B_{\rm Proposal},B_{\rm History}$,
verification must establish
\[
\begin{array}{rcl}
 |\text{original Ready dependencies}+\mathit{OC}|&\le&B_{\rm Ready},\\
 |\text{original NEWVIEW dependencies}+\mathit{OC}+\text{advertisement}|&\le&B_{\rm Newview},\\
 |\text{complete amended REPORT}|&\le&B_{\rm Report},\\
 |\text{original proposal dependencies}+W_R|&\le&B_{\rm Proposal},\\
 |\text{original history dependencies}+\text{required sidecars}|&\le&B_{\rm History}.
\end{array}
\]
The specified handlers, durable attachment, multicast/collection dispatch
and their actual aggregate load must also lie in the original admitted
processing/load profile. The source supplies no numerical codec,
capacity or benchmark proving those inequalities. Finiteness does not
create headroom in a saturated envelope, and separate admissible packets
do not automatically finish in the same $\Delta$. Pre-Ready fragmentation
establishes eventual preparation, not the timed first-Ready or REPORT
bound. The entry-skew and REPORT arithmetic apply at the unchanged
complete-handler abstraction when these explicitly listed dependencies
are admitted. They introduce no new availability or latency premise and
make no unconditional deployment-timing or native-import claim.
